% ============================================================================
%  CHƯƠNG 4a. DỮ LIỆU
% ============================================================================
\clearpage
\section{Dữ liệu}
\label{sec:data}

\subsection{Các bộ dữ liệu}

\begin{table}[H]
\centering
\caption{Vai trò của từng bộ dữ liệu.}
\label{tab:datasets}
\small
\begin{tabularx}{\textwidth}{@{}p{2.9cm} p{3.2cm} L@{}}
\toprule
\textbf{Vai trò} & \textbf{Bộ dữ liệu} & \textbf{Nhãn dùng và ghi chú} \\
\midrule
Huấn luyện, validation, test trong phân phối & VinBigData / VinDr-CXR~\cite{nguyen2022vindr} & 15.000 ảnh huấn
luyện, box theo từng bác sĩ (3 bác sĩ/ảnh). Dương: có ít nhất một box ``Nodule/Mass''; âm: mọi ảnh còn lại, gồm
cả ảnh ``No finding'' và ảnh có bệnh khác (âm khó). \\
\addlinespace
Test ngoài phân phối 1 & NIH ChestX-ray14~\cite{wang2017chestxray8} & Dương: các ảnh có box Nodule hoặc Mass
trong \code{BBox\_List\_2017.csv} (tập con khoảng 1.000 ảnh có box cho 8 bệnh). Âm: ảnh ``No Finding'' của
bệnh nhân không có ảnh dương, lấy ngẫu nhiên gấp 3 lần số ảnh dương (\code{--neg-ratio 3}). Nhãn ``No Finding''
rút từ báo cáo nên có nhiễu. \\
\addlinespace
Test ngoài phân phối 2 (độ mờ) & JSRT~\cite{shiraishi2000jsrt} & 247 ảnh, 154 có nốt kèm vị trí, kích thước
(mm), độ mờ 1--5, lành/ác. Cho phép so trực tiếp với~\cite{nishio2026zeroshot}. \\
\addlinespace
Chỉ tham khảo & NODE21~\cite{sogancioglu2024node21} & Trùng nguồn với JSRT và NIH nên không dùng làm test, để tránh rò rỉ. \\
\bottomrule
\end{tabularx}
\end{table}

Trước khi dùng, nhóm phải kiểm tra điều khoản của cuộc thi VinBigData trên Kaggle, giấy phép NIH và JSRT, và
chấp nhận điều khoản Health AI Developer Foundations để tải MedGemma. Ảnh không được phân phối lại; đồ án chỉ bàn
giao nhãn và dự đoán.

\subsection{Gán nhãn và hợp nhất box}

Mỗi ảnh VinBigData được 3 bác sĩ gán nhãn độc lập. Ảnh là dương nếu có ít nhất một box Nodule/Mass. Box của các
bác sĩ được hợp nhất bằng WBF~\cite{solovyev2021wbf} để có một box tham chiếu cho mỗi tổn thương. Nốt và khối
được gộp thành một lớp dương; vì khối lớn trên 3\,cm dễ phát hiện hơn nhiều, kết quả được phân tầng theo kích thước
để tránh việc khối lớn che lấp hiệu năng trên nốt nhỏ.

\subsection{Chia dữ liệu}

VinBigData được chia theo \code{image\_id}, tỷ lệ 70/10/20, phân tầng theo nhãn Nodule/Mass, với seed cố định
(\code{prepare\_vinbig.py --seed 0}). Bản VinBigData trên Kaggle không kèm mã bệnh nhân, nên không chia được theo
bệnh nhân; nếu một bệnh nhân có nhiều ảnh, các ảnh đó có thể rơi vào cả train và test (một hạn chế, xem
Mục~\ref{sec:conclusion}). Ngưỡng và lựa chọn cấu hình chỉ dùng validation. NIH và JSRT chỉ dùng làm test.
Hình~\ref{fig:split} tóm tắt cách chia, vai trò của từng tập và định nghĩa ca dương.

\begin{figure}[H]
\centering
\begin{tikzpicture}[x=1cm,y=1cm, fb/.append style={font=\scriptsize, minimum height=0.8cm}]
  % ---- thanh chia VinBigData ----
  \node[font=\footnotesize\bfseries, anchor=south west] at (0,0.8) {VinBigData / VinDr-CXR: 15.000 ảnh, $\approx$2,7\% dương};
  \draw[fill=uitblue!25, draw=uitblue] (0,0) rectangle (6.3,0.7);
  \draw[fill=warnorange!25, draw=warnorange] (6.3,0) rectangle (7.2,0.7);
  \draw[fill=badred!18, draw=badred] (7.2,0) rectangle (9.0,0.7);
  \node[font=\footnotesize] at (3.15,0.35) {train 70\%};
  \node[font=\scriptsize] at (6.75,0.35) {val};
  \node[font=\footnotesize] at (8.1,0.35) {test 20\%};
  \node[fnote, anchor=north] at (3.15,-0.05) {$\approx$287 dương (ước tính)};
  \node[fnote, anchor=north] at (6.75,-0.05) {$\approx$41};
  \node[fnote, anchor=north] at (8.1,-0.05) {$\approx$82};
  \node[fnote, anchor=north west] at (0,-0.45) {chia theo \code{image\_id}, phân tầng theo nhãn, seed 0};
  % ---- định nghĩa nhãn ----
  \node[fb, text width=5.7cm, align=left, anchor=north west] (def) at (9.5,1.3) {%
    \textbf{Định nghĩa nhãn (mức ảnh)}\\[1pt]
    \textcolor{badred}{\textbf{Dương}}: $\ge 1$ bác sĩ vẽ box \emph{Nodule/Mass} (nốt và khối gộp một lớp); box hợp nhất bằng WBF.\\[1pt]
    \textcolor{okgreen}{\textbf{Âm}}: mọi ảnh còn lại, kể cả ảnh có bệnh khác.};
  % ---- vai trò ----
  \node[fblue, text width=3.3cm, anchor=north] (tr) at (3.15,-1.35) {\textbf{Huấn luyện LoRA, B1, B2}\\âm:dương $\le$ 4:1 (C3: 1:1)};
  \node[forange, text width=3.3cm, anchor=north] (va) at (7.0,-1.35) {\textbf{Chọn ngưỡng và cấu hình}\\$t_{90}$, $t_{95}$, \code{val\_sens90}};
  \node[fred, text width=3.5cm, anchor=north] (te) at (11.4,-1.35) {\textbf{Đánh giá một lần}\\\code{test}, \code{test\_nih}, \code{test\_jsrt}};
  \draw[farr] (3.15,-0.85) -- (tr.north);
  \draw[farr] (6.75,-0.4) -- (6.75,-1.35);
  \draw[farr] (8.1,-0.4) -- (8.1,-1.05) -| ([xshift=-0.9cm]te.north);
  % ---- tập ngoài ----
  \node[fb, text width=4.4cm, anchor=north west] (nih) at (0,-3.3) {\textbf{NIH ChestX-ray14} (\code{test\_nih})\\dương: ảnh có box Nodule/Mass\\âm: ``No Finding'' $\times 3$, khác bệnh nhân};
  \node[fb, text width=4.4cm, anchor=north west] (js) at (4.95,-3.3) {\textbf{JSRT} (\code{test\_jsrt})\\154 ảnh có nốt (tâm, mm, độ mờ 1--5)\\93 ảnh không có nốt};
  \node[fnote, anchor=north west, text width=5.2cm, align=left] at (10.1,-3.35) {Hai tập ngoài phân phối không tham gia huấn
    luyện hay chọn ngưỡng; ngưỡng $t_{90}$ lấy từ VinBigData val được áp nguyên.};
  \draw[farr] (nih.north) -- ++(0,0.3) -| ([xshift=0.2cm]te.south);
  \draw[farr] (js.north) -- ++(0,0.3) -| ([xshift=0.2cm]te.south);
\end{tikzpicture}
\caption{Chia dữ liệu và định nghĩa nhãn. VinBigData chia 70/10/20 theo ảnh; validation chỉ dùng để chọn ngưỡng
và cấu hình; ba tập test (một trong phân phối, hai ngoài phân phối) chỉ được đánh giá một lần. Số ca dương là ước
tính từ~\cite{nguyen2022vindr}; số thật lấy từ \code{e0\_data.md} sau E0.}
\label{fig:split}
\end{figure}

\begin{warnbox}
\textbf{Rủi ro cỡ mẫu.} Tập huấn luyện của VinDr-CXR chỉ có 410/15.000 ảnh Nodule/Mass
(2,73\%)~\cite{nguyen2022vindr}. Chia 70/10/20 còn khoảng 287 / 41 / 82 ảnh dương (ước tính), nên khoảng tin cậy
của sensitivity trên test sẽ rộng (với 82 ca dương và sensitivity 0,60, nửa độ rộng CI 95\% xấp xỉ $\pm 0{,}11$).
Nếu truy cập được, đồ án ưu tiên dùng tập test chính thức của VinDr-CXR (3.000 ảnh, 176 ảnh Nodule/Mass, nhãn
đồng thuận 5 bác sĩ).
\end{warnbox}

\subsection{Tiền xử lý}

\begin{itemize}
  \item DICOM (VinBigData), PNG (NIH) và ảnh thô 12-bit (JSRT) được chuẩn hoá cường độ, chuyển thành PNG xám
        896$\times$896; hệ số quy đổi toạ độ được lưu lại để đổi box về ảnh gốc.
  \item Qwen2.5-VL được ép cùng số pixel ($896 \times 896$) để hai VLM nhận cùng đầu vào.
  \item Kích thước nốt quy đổi theo diện tích box tương đối; riêng JSRT có kích thước tính bằng mm.
\end{itemize}

\subsection{Thống kê mô tả}

Bảng~\ref{tab:data-stats} được điền từ \code{outputs/e0\_data.md} (sinh bởi \code{data\_stats.py} ở E0).

\begin{table}[H]
\centering
\caption{Thống kê dữ liệu sau tiền xử lý (điền sau E0).}
\label{tab:data-stats}
\small
\begin{tabular}{@{}lrrrr@{}}
\toprule
\textbf{Tập} & \textbf{Số ảnh} & \textbf{Dương} & \textbf{Tỷ lệ dương} & \textbf{Diện tích box trung vị} \\
\midrule
VinBigData train & \textcolor{badred}{[--]} & \textcolor{badred}{[--]} & \textcolor{badred}{[--]} & \textcolor{badred}{[--]} \\
VinBigData val   & \textcolor{badred}{[--]} & \textcolor{badred}{[--]} & \textcolor{badred}{[--]} & \textcolor{badred}{[--]} \\
VinBigData test  & \textcolor{badred}{[--]} & \textcolor{badred}{[--]} & \textcolor{badred}{[--]} & \textcolor{badred}{[--]} \\
NIH ChestX-ray14 (test) & \textcolor{badred}{[--]} & \textcolor{badred}{[--]} & \textcolor{badred}{[--]} & \textcolor{badred}{[--]} \\
JSRT (test)      & 247 & 154 & 62\% & \textcolor{badred}{[--]} \\
\bottomrule
\end{tabular}
\end{table}
